\section{Measurement-accessible tangent geometry}

\subsection{Score-space representation}
Consider a parameterized state measured by a fixed commuting measurement with outcome probabilities $p_a(\bm\theta)$ on a known support $\Omega$. For outcome $a\in\Omega$, define the classical score
\begin{equation}
 s_{ai}=\partial_i\log p_a,
\end{equation}
and the probability-weighted score matrix
\begin{equation}
 S_{ai}=\sqrt{p_a}\,s_{ai}.
\end{equation}
By the definition of classical Fisher information,
\begin{equation}
 G_{\rm C}=\sum_a p_a s_a s_a^{\mathsf T}=S^{\mathsf T}S.
\end{equation}
Moreover, differentiating $\sum_a p_a=1$ gives $\sum_a p_a s_{ai}=\sum_a\partial_i p_a=0$, so every score vector lies in the centered outcome subspace. For any fixed quantum measurement, the associated CFIM is bounded by the QFIM in the same convention, $G_{\rm C}\preceq G_{\rm Q}$ \cite{Braunstein1994}.

Now retain $r$ commuting readout functions $f_1,\ldots,f_r$ of the same outcomes. Let $F\in\mathbb R^{|\Omega|\times r}$ contain these functions and let
\begin{equation}
 \widetilde F_{a\alpha}=F_{a\alpha}-\sum_b p_bF_{b\alpha}
\end{equation}
be their centered version. Introduce
\begin{equation}
 B=D_p^{1/2}\widetilde F,
 \qquad
 \Sigma=B^{\mathsf T}B,
\end{equation}
where $D_p=\operatorname{diag}(p)$. Direct differentiation of the readout expectations gives
\begin{align}
J_{\alpha i}
&=\partial_i\sum_a p_a F_{a\alpha}
 =\sum_a p_a s_{ai}F_{a\alpha}\\
&=\sum_a p_a s_{ai}\widetilde F_{a\alpha},
\end{align}
where the last equality uses the centered-score identity above. Hence
\begin{equation}
 J=\partial_{\bm\theta}\langle F\rangle=B^{\mathsf T}S.
\end{equation}
AQNG uses the covariance-normalized pullback
\begin{equation}
 \boxed{\Gacc=J^{\mathsf T}\Sigma^{+}J.}
 \label{eq:gacc}
\end{equation}

\subsection{Exact projection identity}
Equation~\eqref{eq:gacc} has a direct geometric interpretation. Substituting $J=B^{\mathsf T}S$ and $\Sigma=B^{\mathsf T}B$ gives
\begin{align}
 \Gacc
 &=S^{\mathsf T}B(B^{\mathsf T}B)^{+}B^{\mathsf T}S\\
 &=S^{\mathsf T}P_BS,
 \label{eq:projection_identity}
\end{align}
where $P_B=B(B^{\mathsf T}B)^+B^{\mathsf T}$ is the orthogonal projector onto the probability-weighted centered readout span $\operatorname{col}(B)$. Since $0\preceq P_B\preceq I$,
\begin{equation}
 0\preceq \Gacc\preceq G_{\rm C}\preceq G_{\rm Q},
\end{equation}
with the last inequality supplied by the quantum Fisher bound above \cite{Braunstein1994}. Thus, for commuting retained features, the accessible metric is exactly the Gram matrix of the projected measurement score. It is not defined as a generic low-rank approximation to $G_{\rm Q}$; it is the metric induced by a specified measurement interface. Figure~\ref{fig:accessible_geometry} summarizes the construction.

\begin{figure}[t]
\centering
\begin{tikzpicture}[
  node distance=5mm,
  box/.style={draw, rounded corners, align=center, inner sep=4pt, font=\scriptsize},
  >=Latex
]
\node[box] (score) {measurement score\\$S$};
\node[box, right=of score] (proj) {retained span\\$P_B$};
\node[box, right=of proj] (metric) {accessible metric\\$\Gacc=S^{\mathsf T}P_BS$};
\node[box, below=of metric] (step) {preconditioned step\\$(\Gacc+\lambda I)d=g$};
\draw[->] (score) -- (proj);
\draw[->] (proj) -- (metric);
\draw[->] (metric) -- (step);
\node[font=\scriptsize, below=6mm of score] (full) {$G_{\rm C}=S^{\mathsf T}S$};
\draw[->] (full) -- (score);
\end{tikzpicture}
\caption{Accessible natural-gradient construction. The retained readout span projects the measurement score before the pullback metric is formed; the resulting metric preconditions the task gradient.}
\label{fig:accessible_geometry}
\end{figure}

The covariance normalization removes arbitrary basis scaling inside a fixed retained span. Indeed, replacing $B$ by $BM$ for any nonsingular $M$ leaves $\operatorname{col}(B)$ unchanged, hence leaves the orthogonal projector $P_B$ and therefore $\Gacc$ unchanged by Eq.~\eqref{eq:projection_identity}.

\subsection{Retention and orientation}
To diagnose which score directions a rank-$r$ readout captures, we use the normalized tangent-score covariance introduced in Refs.~\cite{AitHaddou2026ReadoutRank,AitHaddou2026SpectralGeometry}, $C\succeq0$ with $\Tr C=1$. If $P$ projects onto the retained score subspace, the retained mass and rank-normalized retention are
\begin{equation}
 R=\Tr(PC),
 \qquad
 \rho=\frac{R}{r/N},
 \label{eq:rho}
\end{equation}
where $N$ is the centered score-space dimension. Under a uniformly random rank-$r$ subspace, $\mathbb E\rho=1$ \cite{AitHaddou2026ReadoutRank,AitHaddou2026SpectralGeometry}. Values $\rho\gg1$ therefore diagnose orientation beyond the rank baseline. The corresponding random-subspace variance is controlled by the spectrum of $C$, which separates readout rank, tangent anisotropy, and relative orientation \cite{AitHaddou2026SpectralGeometry}.

\subsection{Scope}
Equation~\eqref{eq:projection_identity} is an operational statement only when the retained functions are jointly estimable from the same measurement outcome record. In this work that condition is enforced by using diagonal functions of computational-basis outcomes. If one instead wishes to combine incompatible observables, the measurement protocol and its joint classical record must first be specified; only then can an analogous score-space construction be defined.
